\documentclass[11pt,a4paper]{article}
\usepackage[utf8]{inputenc}
\usepackage[T1]{fontenc}
\usepackage[catalan]{babel}
\usepackage{lmodern}
\usepackage[sfdefault,lining]{FiraSans}
\usepackage[a4paper,top=2.3cm,bottom=2.2cm,left=1.8cm,right=1.8cm,headheight=1.3cm,headsep=0.4cm,footskip=1cm]{geometry}
\usepackage{amsmath,amssymb}
\usepackage{icomma}
\usepackage{graphicx}
\usepackage[table]{xcolor}
\usepackage{tikz}
\usepackage{fontawesome5}
\usepackage{ccicons}
\usepackage{fancyhdr}
\usepackage{lastpage}
\usepackage{tasks}
\usepackage{needspace}
\usepackage[most]{tcolorbox}
\usepackage{array}
\usepackage[hidelinks,pdftitle={Solucions de la fitxa de la unitat 5: Equacions. 3r d'ESO},pdfsubject={Matemàtiques 3r ESO},pdfkeywords={equacions, identitats, equacions de primer grau, denominadors, equacions de segon grau, discriminant, equacions de tercer grau, regla de Ruffini, problemes}]{hyperref}

\definecolor{navy}{HTML}{1B2748}
\definecolor{accent}{HTML}{E8603A}
\definecolor{teal}{HTML}{1F9E90}
\definecolor{soft}{HTML}{EEF1F8}
\definecolor{muted}{HTML}{5D6685}

\setlength{\parindent}{0pt}
\setlength{\parskip}{3pt}
\settasks{label=\alph*), label-format=\bfseries\color{navy}, label-width=1.3em, item-indent=1.8em, column-sep=0.6em, after-item-skip=5pt, before-skip=3pt, after-skip=4pt}

% ---------- capçalera i peu ----------
\newcommand{\llicencia}{\href{https://creativecommons.org/licenses/by-nc-sa/4.0/deed.ca}{\ccbyncsa}\ \ Llicència Creative Commons Reconeixement-NoComercial-CompartirIgual 4.0 Internacional (CC BY-NC-SA 4.0)}
\pagestyle{fancy}
\fancyhf{}
\fancyhead[L]{\raisebox{-0.35\height}{\includegraphics[height=0.95cm]{logo-gva.png}}}
\fancyhead[C]{\small\color{muted}Solucions de la fitxa · Unitat 5: Equacions · 3r d'ESO}
\fancyhead[R]{\raisebox{-0.35\height}{\includegraphics[height=0.95cm]{logo-vermellar.png}}}
\fancyfoot[L]{\scriptsize\color{muted}\llicencia}
\fancyfoot[R]{\small\color{muted}\thepage}
\renewcommand{\headrulewidth}{0.4pt}
\renewcommand{\headrule}{\color{muted!40}\hrule height \headrulewidth}
\fancypagestyle{primera}{\fancyhf{}\renewcommand{\headrulewidth}{0pt}
  \fancyfoot[L]{\scriptsize\color{muted}\llicencia}
  \fancyfoot[R]{\small\color{muted}\thepage}}

% ---------- nivells i exercicis ----------
\newcommand{\estrelles}[1]{{\color{accent}\ifcase#1\or\faStar\faStar[regular]\faStar[regular]\or\faStar\faStar\faStar[regular]\or\faStar\faStar\faStar\fi}}
\newcounter{exer}
\newcommand{\ex}[1]{\par\needspace{4\baselineskip}\medskip\stepcounter{exer}\noindent{\color{navy}\bfseries\theexer.}\ \ {\small\estrelles{#1}}\ \ }

\newcommand{\seccio}[2]{\par\needspace{10\baselineskip}\bigskip
  \begin{tcolorbox}[colback=navy,colframe=navy,coltext=white,arc=2mm,boxrule=0pt,left=3mm,right=3mm,top=1.2mm,bottom=1.2mm]
  {\large\bfseries #1\quad #2}
  \end{tcolorbox}\vspace{-2pt}}

\newcolumntype{C}[1]{>{\centering\arraybackslash}p{#1}}

% quadre de recordatori (només a la fitxa d'exercicis)
\newcommand{\recorda}[1]{\par\smallskip\begin{tcolorbox}[colback=accent!7,colframe=accent!55,arc=2mm,boxrule=0.5pt,left=2.5mm,right=2.5mm,top=1mm,bottom=1mm]\small{\color{accent}\faLightbulb[regular]}\ \textbf{Recorda:}\ #1\end{tcolorbox}}
\newcommand{\res}[1]{\boxed{#1}}
\newcommand{\rest}[1]{\fbox{#1}}

\def\fmneed{14}
% =====================================================================
% Capçalera FentMates de les fitxes: un color per a cada curs.
% Els colors són els de la pàgina inicial (LEVELS a assets/js/materials.js);
% el logo ve de _eines/logo/pdf/ (python3 _eines/logo/genera.py).
%
% Ús, al final del preàmbul de la unitat:
%   % =====================================================================
% Capçalera FentMates de les fitxes: un color per a cada curs.
% Els colors són els de la pàgina inicial (LEVELS a assets/js/materials.js);
% el logo ve de _eines/logo/pdf/ (python3 _eines/logo/genera.py).
%
% Ús, al final del preàmbul de la unitat:
%   % =====================================================================
% Capçalera FentMates de les fitxes: un color per a cada curs.
% Els colors són els de la pàgina inicial (LEVELS a assets/js/materials.js);
% el logo ve de _eines/logo/pdf/ (python3 _eines/logo/genera.py).
%
% Ús, al final del preàmbul de la unitat:
%   \input{../comu/fentmates.tex}
%   \fmcurs{1eso}
% i, al començament del document:
%   \capcalera{Matemàtiques}{1r ESO}{Fitxa d'exercicis · Unitat 2}{Divisibilitat}
% Opcional abans de \input: \def\fmneed{14} (espai mínim abans de cada secció).
% =====================================================================
\makeatletter
\def\fm@logodir{../../logo/pdf}
\providecommand{\fmneed}{10}
\def\fm@cursnom{}

% clau del curs, color i color del text damunt del color
\newcommand{\fm@curs}[3]{\@namedef{fm@color@#1}{#2}\@namedef{fm@text@#1}{#3}}
\fm@curs{1eso}{E6B84A}{navy}
\fm@curs{2eso}{3B82F6}{white}
\fm@curs{3eso}{FF7A50}{white}
\fm@curs{4eso-a}{A78BFA}{white}
\fm@curs{4eso-b}{6D28D9}{white}
\fm@curs{1batx-mi}{9B1C1C}{white}
\fm@curs{1batx-ccss}{EF6B6B}{white}
\fm@curs{2batx-mii}{1E6B47}{white}
\fm@curs{2batx-ccss}{2BB3A3}{white}

\newcommand{\fmcurs}[1]{%
  \@ifundefined{fm@color@#1}{\PackageError{fentmates}{Curs desconegut: #1}{}}{}%
  \edef\fm@tmp{\noexpand\definecolor{curs}{HTML}{\@nameuse{fm@color@#1}}%
               \noexpand\colorlet{curstext}{\@nameuse{fm@text@#1}}}\fm@tmp
  \colorlet{cursfosc}{curs!72!black}%
  \def\fm@logo{\fm@logodir/vertical-#1.pdf}%
  \def\fm@marca{\fm@logodir/marca-#1.pdf}}

% ---------- capçalera de la primera pàgina ----------
% #1 matèria (petit), #2 curs, #3 tipus de document i unitat, #4 títol
\newcommand{\capcalera}[4]{%
  \gdef\fm@cursnom{#2}%
  \par\noindent
  \begin{tcolorbox}[colback=curs,colframe=curs,boxrule=0pt,arc=2.2mm,width=3.5cm,height=2.15cm,
      valign=center,left=3mm,right=2mm,top=1mm,bottom=1mm,nobeforeafter,box align=center]
    \color{curstext}{\footnotesize #1}\par\vspace{1pt}{\fontsize{19}{21}\selectfont\bfseries #2\par}%
  \end{tcolorbox}\hspace{5mm}%
  \begin{minipage}[c]{\dimexpr\linewidth-3.5cm-5mm-2.9cm\relax}
    {\color{cursfosc}\small\bfseries\MakeUppercase{#3}}\par\vspace{2pt}
    {\color{navy}\fontsize{20}{23}\selectfont\bfseries #4\par}
  \end{minipage}\hfill
  \begin{minipage}[c]{2.6cm}\raggedleft\includegraphics[width=2.6cm]{\fm@logo}\end{minipage}\par
  \vspace{2pt}}

% ---------- barres de secció del color del curs ----------
\renewcommand{\seccio}[2]{\par\needspace{\fmneed\baselineskip}\bigskip
  \begin{tcolorbox}[colback=curs,colframe=curs,coltext=curstext,arc=2mm,boxrule=0pt,left=3mm,right=3mm,top=1.2mm,bottom=1.2mm]
  {\large\bfseries #1\quad #2}
  \end{tcolorbox}\vspace{-2pt}}

% ---------- capçalera i peu de les pàgines ----------
% primera pàgina: logos institucionals al peu
\fancypagestyle{primera}{\fancyhf{}\renewcommand{\headrulewidth}{0pt}
  \fancyfoot[L]{\raisebox{-0.32\height}{\includegraphics[height=0.78cm]{logo-gva.png}}\hspace{5mm}%
    \raisebox{-0.32\height}{\includegraphics[height=0.78cm]{logo-vermellar.png}}}
  \fancyfoot[R]{\scriptsize\color{muted}\href{https://creativecommons.org/licenses/by-nc-sa/4.0/deed.ca}{\ccbyncsa}\ \ CC BY-NC-SA 4.0 · fentmates.org\qquad\small\thepage}}
% la resta: el símbol a l'esquerra i el curs a la dreta, amb el títol al centre com abans
\fancyhead[L]{\raisebox{-0.3\height}{\includegraphics[height=0.62cm]{\fm@marca}}}
\fancyhead[R]{\small\bfseries\color{cursfosc}\fm@cursnom}
\renewcommand{\headrule}{\color{curs!45}\hrule height \headrulewidth}
\makeatother

%   \fmcurs{1eso}
% i, al començament del document:
%   \capcalera{Matemàtiques}{1r ESO}{Fitxa d'exercicis · Unitat 2}{Divisibilitat}
% Opcional abans de \input: \def\fmneed{14} (espai mínim abans de cada secció).
% =====================================================================
\makeatletter
\def\fm@logodir{../../logo/pdf}
\providecommand{\fmneed}{10}
\def\fm@cursnom{}

% clau del curs, color i color del text damunt del color
\newcommand{\fm@curs}[3]{\@namedef{fm@color@#1}{#2}\@namedef{fm@text@#1}{#3}}
\fm@curs{1eso}{E6B84A}{navy}
\fm@curs{2eso}{3B82F6}{white}
\fm@curs{3eso}{FF7A50}{white}
\fm@curs{4eso-a}{A78BFA}{white}
\fm@curs{4eso-b}{6D28D9}{white}
\fm@curs{1batx-mi}{9B1C1C}{white}
\fm@curs{1batx-ccss}{EF6B6B}{white}
\fm@curs{2batx-mii}{1E6B47}{white}
\fm@curs{2batx-ccss}{2BB3A3}{white}

\newcommand{\fmcurs}[1]{%
  \@ifundefined{fm@color@#1}{\PackageError{fentmates}{Curs desconegut: #1}{}}{}%
  \edef\fm@tmp{\noexpand\definecolor{curs}{HTML}{\@nameuse{fm@color@#1}}%
               \noexpand\colorlet{curstext}{\@nameuse{fm@text@#1}}}\fm@tmp
  \colorlet{cursfosc}{curs!72!black}%
  \def\fm@logo{\fm@logodir/vertical-#1.pdf}%
  \def\fm@marca{\fm@logodir/marca-#1.pdf}}

% ---------- capçalera de la primera pàgina ----------
% #1 matèria (petit), #2 curs, #3 tipus de document i unitat, #4 títol
\newcommand{\capcalera}[4]{%
  \gdef\fm@cursnom{#2}%
  \par\noindent
  \begin{tcolorbox}[colback=curs,colframe=curs,boxrule=0pt,arc=2.2mm,width=3.5cm,height=2.15cm,
      valign=center,left=3mm,right=2mm,top=1mm,bottom=1mm,nobeforeafter,box align=center]
    \color{curstext}{\footnotesize #1}\par\vspace{1pt}{\fontsize{19}{21}\selectfont\bfseries #2\par}%
  \end{tcolorbox}\hspace{5mm}%
  \begin{minipage}[c]{\dimexpr\linewidth-3.5cm-5mm-2.9cm\relax}
    {\color{cursfosc}\small\bfseries\MakeUppercase{#3}}\par\vspace{2pt}
    {\color{navy}\fontsize{20}{23}\selectfont\bfseries #4\par}
  \end{minipage}\hfill
  \begin{minipage}[c]{2.6cm}\raggedleft\includegraphics[width=2.6cm]{\fm@logo}\end{minipage}\par
  \vspace{2pt}}

% ---------- barres de secció del color del curs ----------
\renewcommand{\seccio}[2]{\par\needspace{\fmneed\baselineskip}\bigskip
  \begin{tcolorbox}[colback=curs,colframe=curs,coltext=curstext,arc=2mm,boxrule=0pt,left=3mm,right=3mm,top=1.2mm,bottom=1.2mm]
  {\large\bfseries #1\quad #2}
  \end{tcolorbox}\vspace{-2pt}}

% ---------- capçalera i peu de les pàgines ----------
% primera pàgina: logos institucionals al peu
\fancypagestyle{primera}{\fancyhf{}\renewcommand{\headrulewidth}{0pt}
  \fancyfoot[L]{\raisebox{-0.32\height}{\includegraphics[height=0.78cm]{logo-gva.png}}\hspace{5mm}%
    \raisebox{-0.32\height}{\includegraphics[height=0.78cm]{logo-vermellar.png}}}
  \fancyfoot[R]{\scriptsize\color{muted}\href{https://creativecommons.org/licenses/by-nc-sa/4.0/deed.ca}{\ccbyncsa}\ \ CC BY-NC-SA 4.0 · fentmates.org\qquad\small\thepage}}
% la resta: el símbol a l'esquerra i el curs a la dreta, amb el títol al centre com abans
\fancyhead[L]{\raisebox{-0.3\height}{\includegraphics[height=0.62cm]{\fm@marca}}}
\fancyhead[R]{\small\bfseries\color{cursfosc}\fm@cursnom}
\renewcommand{\headrule}{\color{curs!45}\hrule height \headrulewidth}
\makeatother

%   \fmcurs{1eso}
% i, al començament del document:
%   \capcalera{Matemàtiques}{1r ESO}{Fitxa d'exercicis · Unitat 2}{Divisibilitat}
% Opcional abans de \input: \def\fmneed{14} (espai mínim abans de cada secció).
% =====================================================================
\makeatletter
\def\fm@logodir{../../logo/pdf}
\providecommand{\fmneed}{10}
\def\fm@cursnom{}

% clau del curs, color i color del text damunt del color
\newcommand{\fm@curs}[3]{\@namedef{fm@color@#1}{#2}\@namedef{fm@text@#1}{#3}}
\fm@curs{1eso}{E6B84A}{navy}
\fm@curs{2eso}{3B82F6}{white}
\fm@curs{3eso}{FF7A50}{white}
\fm@curs{4eso-a}{A78BFA}{white}
\fm@curs{4eso-b}{6D28D9}{white}
\fm@curs{1batx-mi}{9B1C1C}{white}
\fm@curs{1batx-ccss}{EF6B6B}{white}
\fm@curs{2batx-mii}{1E6B47}{white}
\fm@curs{2batx-ccss}{2BB3A3}{white}

\newcommand{\fmcurs}[1]{%
  \@ifundefined{fm@color@#1}{\PackageError{fentmates}{Curs desconegut: #1}{}}{}%
  \edef\fm@tmp{\noexpand\definecolor{curs}{HTML}{\@nameuse{fm@color@#1}}%
               \noexpand\colorlet{curstext}{\@nameuse{fm@text@#1}}}\fm@tmp
  \colorlet{cursfosc}{curs!72!black}%
  \def\fm@logo{\fm@logodir/vertical-#1.pdf}%
  \def\fm@marca{\fm@logodir/marca-#1.pdf}}

% ---------- capçalera de la primera pàgina ----------
% #1 matèria (petit), #2 curs, #3 tipus de document i unitat, #4 títol
\newcommand{\capcalera}[4]{%
  \gdef\fm@cursnom{#2}%
  \par\noindent
  \begin{tcolorbox}[colback=curs,colframe=curs,boxrule=0pt,arc=2.2mm,width=3.5cm,height=2.15cm,
      valign=center,left=3mm,right=2mm,top=1mm,bottom=1mm,nobeforeafter,box align=center]
    \color{curstext}{\footnotesize #1}\par\vspace{1pt}{\fontsize{19}{21}\selectfont\bfseries #2\par}%
  \end{tcolorbox}\hspace{5mm}%
  \begin{minipage}[c]{\dimexpr\linewidth-3.5cm-5mm-2.9cm\relax}
    {\color{cursfosc}\small\bfseries\MakeUppercase{#3}}\par\vspace{2pt}
    {\color{navy}\fontsize{20}{23}\selectfont\bfseries #4\par}
  \end{minipage}\hfill
  \begin{minipage}[c]{2.6cm}\raggedleft\includegraphics[width=2.6cm]{\fm@logo}\end{minipage}\par
  \vspace{2pt}}

% ---------- barres de secció del color del curs ----------
\renewcommand{\seccio}[2]{\par\needspace{\fmneed\baselineskip}\bigskip
  \begin{tcolorbox}[colback=curs,colframe=curs,coltext=curstext,arc=2mm,boxrule=0pt,left=3mm,right=3mm,top=1.2mm,bottom=1.2mm]
  {\large\bfseries #1\quad #2}
  \end{tcolorbox}\vspace{-2pt}}

% ---------- capçalera i peu de les pàgines ----------
% primera pàgina: logos institucionals al peu
\fancypagestyle{primera}{\fancyhf{}\renewcommand{\headrulewidth}{0pt}
  \fancyfoot[L]{\raisebox{-0.32\height}{\includegraphics[height=0.78cm]{logo-gva.png}}\hspace{5mm}%
    \raisebox{-0.32\height}{\includegraphics[height=0.78cm]{logo-vermellar.png}}}
  \fancyfoot[R]{\scriptsize\color{muted}\href{https://creativecommons.org/licenses/by-nc-sa/4.0/deed.ca}{\ccbyncsa}\ \ CC BY-NC-SA 4.0 · fentmates.org\qquad\small\thepage}}
% la resta: el símbol a l'esquerra i el curs a la dreta, amb el títol al centre com abans
\fancyhead[L]{\raisebox{-0.3\height}{\includegraphics[height=0.62cm]{\fm@marca}}}
\fancyhead[R]{\small\bfseries\color{cursfosc}\fm@cursnom}
\renewcommand{\headrule}{\color{curs!45}\hrule height \headrulewidth}
\makeatother

\fmcurs{3eso}

\begin{document}
\thispagestyle{primera}

\capcalera{Matemàtiques}{3r ESO}{Solucions de la fitxa · Unitat 5}{Equacions}
\vspace{8pt}

\begin{tcolorbox}[colback=soft,colframe=soft,arc=2mm,left=3mm,right=3mm,top=1.5mm,bottom=1.5mm]
\small Intenta resoldre cada exercici abans de mirar-ne la solució. Si no t'ix el mateix resultat, repassa els passos o consulta els apunts interactius de la unitat: hi trobaràs calculadores que resolen pas a pas qualsevol equació.
\end{tcolorbox}

% =====================================================
\seccio{1}{Equacions i identitats}

\ex{1} \begin{tasks}(2)
  \task \rest{Identitat}
  \task \rest{Equació} ($x=3$)
  \task \rest{Identitat}
  \task \rest{Equació} ($x=\pm1$)
  \task \rest{Identitat}
  \task \rest{Equació} ($x=1$)
  \task \rest{Identitat}
  \task \rest{Equació} sense solució ($-2\neq-1$)
\end{tasks}

\ex{1} \begin{tasks}(2)
  \task $4=4$: \rest{sí}
  \task $5=5$ i $5=5$: \rest{totes dues}
  \task $8\neq9$: \rest{no}
  \task $8-4=4$: \rest{sí}
\end{tasks}

\ex{1} 1r membre: $5x^2-3x+1$; 2n membre: $2x-7$; termes: $5x^2,\ -3x,\ 1,\ 2x,\ -7$; incògnita: $x$; \rest{grau 2}.

\ex{2} \begin{tasks}(2)
  \task $4x=-8$: \rest{grau 1}
  \task $-4x=0$: \rest{grau 1}
  \task $-x=2$: \rest{grau 1}
  \task $x^2-3x-1=0$: \rest{grau 2}
\end{tasks}

\ex{2} \begin{tasks}(1)
  \task Les dues tenen la solució $x=5$: \rest{sí}.
  \task La primera també té la solució $x=-5$: \rest{no}.
\end{tasks}

\ex{2} Per exemple, \rest{$x+2=0$} \ (o $3x+1=-5$\ldots) \ i \ \rest{$x^2-9=0$}.

\ex{3} \begin{tasks}(1)
  \task $18-3k+3=0\Rightarrow\res{k=7}$
  \task $-k-2=k-6\Rightarrow -2k=-4\Rightarrow\res{k=2}$
\end{tasks}


% =====================================================
\seccio{2}{Equacions de primer grau}

\ex{1} \begin{tasks}(2)
  \task $\res{x=5}$
  \task $\res{x=-7}$
  \task $5x=15\ \Rightarrow\ \res{x=3}$
  \task $-2x=-8\ \Rightarrow\ \res{x=4}$
  \task $2x=8\ \Rightarrow\ \res{x=4}$
  \task $-4x=-20\ \Rightarrow\ \res{x=5}$
  \task $-3x=12\ \Rightarrow\ \res{x=-4}$
  \task $-3x=0\ \Rightarrow\ \res{x=0}$
\end{tasks}

\ex{1} \begin{tasks}(1)
  \task $3x-6=12\ \Rightarrow\ 3x=18\ \Rightarrow\ \res{x=6}$
  \task $5x+5=2x-7\ \Rightarrow\ 3x=-12\ \Rightarrow\ \res{x=-4}$
  \task $-x+1=2x-5\ \Rightarrow\ -3x=-6\ \Rightarrow\ \res{x=2}$
  \task $6x-2=4x+8\ \Rightarrow\ 2x=10\ \Rightarrow\ \res{x=5}$
\end{tasks}

\ex{2} \begin{tasks}(1)
  \task $4x-11=5x-13\ \Rightarrow\ -x=-2\ \Rightarrow\ \res{x=2}$
  \task $2x-1=4x-7\ \Rightarrow\ -2x=-6\ \Rightarrow\ \res{x=3}$
  \task $-2x+3=7\ \Rightarrow\ -2x=4\ \Rightarrow\ \res{x=-2}$
  \task $3x-7=2x+2\ \Rightarrow\ \res{x=9}$
\end{tasks}

\ex{2} \begin{tasks}(1)
  \task $2x+6=2x+6\ \Rightarrow\ 0x=0\ \Rightarrow\ \res{\text{infinites solucions}}$
  \task $3x+2=3x+4\ \Rightarrow\ 0x=2\ \Rightarrow\ \res{\text{sense solució}}$
  \task $-2=-2\ \Rightarrow\ 0x=0\ \Rightarrow\ \res{\text{infinites solucions}}$
  \task $x+3=2x+1\ \Rightarrow\ -x=-2\ \Rightarrow\ \res{x=2}$
\end{tasks}

\ex{2} \begin{tasks}(1)
  \task $4x=12\ \Rightarrow\ \res{x=3}$. \ Comprovació: $16-4=12$
  \task $3x=2x+5\ \Rightarrow\ \res{x=5}$. \ Comprovació: $40-25=15=2\cdot5+5$
\end{tasks}

\ex{3} \begin{tasks}(1)
  \task $-4x+2=-7x+8\ \Rightarrow\ 3x=6\ \Rightarrow\ \res{x=2}$
  \task $2x-10=2x+2\ \Rightarrow\ 0x=12\ \Rightarrow\ \res{\text{sense solució}}$
\end{tasks}

\ex{3} $3(5-a)=11\Rightarrow 15-3a=11\Rightarrow\res{a=\tfrac43}$


% =====================================================
\seccio{3}{Equacions de primer grau amb denominadors}

\ex{1} \begin{tasks}(2)
  \task $\text{m.c.m.}=3:\ \ \res{x=12}$
  \task $\text{m.c.m.}=4:\ \ 2x+x=24\ \Rightarrow\ 3x=24\ \Rightarrow\ \res{x=8}$
  \task $\text{m.c.m.}=3:\ \ 3x-x=30\ \Rightarrow\ 2x=30\ \Rightarrow\ \res{x=15}$
  \task $\text{m.c.m.}=5:\ \ 2x-5=15\ \Rightarrow\ 2x=20\ \Rightarrow\ \res{x=10}$
\end{tasks}

\ex{1} \begin{tasks}(2)
  \task $\text{m.c.m.}=3:\ \ (x+2)=12\ \Rightarrow\ x+2=12\ \Rightarrow\ \res{x=10}$
  \task $\text{m.c.m.}=4:\ \ 2(x-1)=(x+3)\ \Rightarrow\ 2x-2=x+3\ \Rightarrow\ \res{x=5}$
  \task $\text{m.c.m.}=4:\ \ (3x-1)=8\ \Rightarrow\ 3x-1=8\ \Rightarrow\ 3x=9\ \Rightarrow\ \res{x=3}$
  \task $\text{m.c.m.}=12:\ \ 2x+12=3x\ \Rightarrow\ -x=-12\ \Rightarrow\ \res{x=12}$
\end{tasks}

\ex{2} \begin{tasks}(1)
  \task $\text{m.c.m.}=6:\ \ 3(x+1)-2(x-1)=12\ \Rightarrow\ x+5=12\ \Rightarrow\ \res{x=7}$
  \task $\text{m.c.m.}=12:\ \ 3(x-2)-2(x-3)=6\ \Rightarrow\ \res{x=6}$
  \task $\text{m.c.m.}=6:\ \ 2(2x+1)-3(x-2)=18\ \Rightarrow\ x+8=18\ \Rightarrow\ \res{x=10}$
  \task $\text{m.c.m.}=10:\ \ 10x-2(x+2)=5(x+6)\ \Rightarrow\ 8x-4=5x+30\ \Rightarrow\ 3x=34\ \Rightarrow\ \res{x=\tfrac{34}{3}}$
\end{tasks}

\ex{2} \begin{tasks}(1)
  \task $\text{m.c.m.}=6:\ \ 2(x-1)+(x+2)=6x-12\ \Rightarrow\ 3x=6x-12\ \Rightarrow\ -3x=-12\ \Rightarrow\ \res{x=4}$
  \task $\text{m.c.m.}=4:\ \ 12-(x-1)=2(x+5)\ \Rightarrow\ -x+13=2x+10\ \Rightarrow\ -3x=-3\ \Rightarrow\ \res{x=1}$
\end{tasks}

\ex{3} \begin{tasks}(1)
  \task $\text{m.c.m.}=12:\ \ 9(x-2)-8(x+1)=2(x-10)\ \Rightarrow\ x-26=2x-20\ \Rightarrow\ -x=6\ \Rightarrow\ \res{x=-6}$
  \task $\text{m.c.m.}=12:\ \ 6x-4(x-1)-3(x+1)=(1-x)\ \Rightarrow\ -x+1=-x+1\ \Rightarrow\ 0x=0\ \Rightarrow\ \res{\text{infinites solucions}}$
  \task $\text{m.c.m.}=4:\ \ 2(x+1)-(x+3)=x-4\ \Rightarrow\ x-1=x-4\ \Rightarrow\ 0x=-3\ \Rightarrow\ \res{\text{sense solució}}$
\end{tasks}

\ex{3} El $-$ davant de la segona fracció afecta tot el numerador: $-3(x+1)=-3x-3$. \ $2x-4-3x-3=6\Rightarrow -x=13\Rightarrow\res{x=-13}$.


% =====================================================
\seccio{4}{Equacions de segon grau incompletes}

\ex{1} \begin{tasks}(1)
  \task $x^2-5x+6=0$: \rest{$a=1,\ b=-5,\ c=6$, completa}
  \task $3x^2-12=0$: \rest{$a=3,\ b=0,\ c=-12$, incompleta}
  \task $x^2-4x=0$: \rest{$a=1,\ b=-4,\ c=0$, incompleta}
  \task $x^2-4=0$: \rest{$a=1,\ b=0,\ c=-4$, incompleta}
\end{tasks}

\ex{1} \begin{tasks}(2)
  \task $x^{2}=36\ \Rightarrow\ \res{x=\pm6}$
  \task $x^{2}=25\ \Rightarrow\ \res{x=\pm5}$
  \task $x^{2}=-9\ \Rightarrow\ \res{\text{sense solució}}$
  \task $x^{2}=\tfrac{9}{4}\ \Rightarrow\ \res{x=\pm\tfrac{3}{2}}$
  \task $x^{2}=9\ \Rightarrow\ \res{x=\pm3}$
  \task $x^{2}=7\ \Rightarrow\ \res{x=\pm\sqrt{7}}$
\end{tasks}

\ex{1} \begin{tasks}(1)
  \task $x\,(x-5)=0\ \Rightarrow\ \res{x=0,\ x=5}$
  \task $x^{2}+4x=0\ \Rightarrow\ x\,(x+4)=0\ \Rightarrow\ \res{x=-4,\ x=0}$
  \task $2x^{2}-7x=0\ \Rightarrow\ x\,(2x-7)=0\ \Rightarrow\ \res{x=0,\ x=\tfrac{7}{2}}$
  \task $x^{2}-6x=0\ \Rightarrow\ x\,(x-6)=0\ \Rightarrow\ \res{x=0,\ x=6}$
\end{tasks}

\ex{2} \begin{tasks}(1)
  \task $x^{2}-9=7\ \Rightarrow\ x^{2}=16\ \Rightarrow\ \res{x=\pm4}$
  \task $x^{2}-2x=-2x+4\ \Rightarrow\ x^{2}=4\ \Rightarrow\ \res{x=\pm2}$
  \task $4x^{2}-4x+1=1\ \Rightarrow\ x^{2}-x=0\ \Rightarrow\ x\,(x-1)=0\ \Rightarrow\ \res{x=0,\ x=1}$
  \task $x^{2}+4x=3x\ \Rightarrow\ x^{2}+x=0\ \Rightarrow\ x\,(x+1)=0\ \Rightarrow\ \res{x=-1,\ x=0}$
\end{tasks}

\ex{2} \begin{tasks}(1)
  \task $\text{m.c.m.}=3:\ \ x^{2}=36\ \Rightarrow\ \res{x=\pm6}$
  \task $\text{m.c.m.}=4:\ \ x^{2}-2x=0\ \Rightarrow\ x\,(x-2)=0\ \Rightarrow\ \res{x=0,\ x=2}$
  \task $\text{m.c.m.}=5:\ \ (x^{2}+1)=10\ \Rightarrow\ x^{2}+1=10\ \Rightarrow\ x^{2}=9\ \Rightarrow\ \res{x=\pm3}$
  \task $\text{m.c.m.}=6:\ \ 3(x^{2}-2)-2(x^{2}+1)=6\ \Rightarrow\ x^{2}-8=6\ \Rightarrow\ x^{2}=14\ \Rightarrow\ \res{x=\pm\sqrt{14}}$
\end{tasks}

\ex{3} \begin{tasks}(1)
  \task $3x^{2}-2x-10=2x^{2}-2x-10\ \Rightarrow\ x^{2}=0\ \Rightarrow\ \res{x=0\ (\text{doble})}$
  \task $\text{m.c.m.}=6:\ \ 3x(x-3)-2x(x+1)=6x\ \Rightarrow\ x^{2}-11x=6x\ \Rightarrow\ x^{2}-17x=0\ \Rightarrow\ x\,(x-17)=0\ \Rightarrow\ \res{x=0,\ x=17}$
\end{tasks}

\ex{3} En dividir entre $x$ perd la solució $x=0$. \ $5x^2-3x=0\Rightarrow x(5x-3)=0\Rightarrow\res{x=0,\ x=\tfrac35}$.

\ex{3} \begin{tasks}(2)
  \task P. ex.\ \rest{$x^2-16=0$}
  \task P. ex.\ \rest{$x^2-3x=0$}
\end{tasks}


% =====================================================
\seccio{5}{Equacions de segon grau completes}

\ex{1} \begin{tasks}(2)
  \task $\Delta=4\ \Rightarrow\ x=\tfrac{6\pm2}{2}\ \Rightarrow\ \res{x=2,\ x=4}$
  \task $\Delta=36\ \Rightarrow\ x=\tfrac{-2\pm6}{2}\ \Rightarrow\ \res{x=-4,\ x=2}$
  \task $\Delta=36\ \Rightarrow\ x=\tfrac{4\pm6}{2}\ \Rightarrow\ \res{x=-1,\ x=5}$
  \task $\Delta=1\ \Rightarrow\ x=\tfrac{-7\pm1}{2}\ \Rightarrow\ \res{x=-4,\ x=-3}$
\end{tasks}

\ex{1} \begin{tasks}(2)
  \task $\Delta=9\ \Rightarrow\ x=\tfrac{5\pm3}{4}\ \Rightarrow\ \res{x=\tfrac{1}{2},\ x=2}$
  \task $\Delta=25\ \Rightarrow\ x=\tfrac{-1\pm5}{6}\ \Rightarrow\ \res{x=-1,\ x=\tfrac{2}{3}}$
  \task $\Delta=25\ \Rightarrow\ x=\tfrac{1\pm5}{12}\ \Rightarrow\ \res{x=-\tfrac{1}{3},\ x=\tfrac{1}{2}}$
  \task $\Delta=64\ \Rightarrow\ x=\tfrac{-4\pm8}{8}\ \Rightarrow\ \res{x=-\tfrac{3}{2},\ x=\tfrac{1}{2}}$
\end{tasks}

\ex{1} \begin{tasks}(2)
  \task $\Delta=0$: \rest{una solució}
  \task $\Delta=-23$: \rest{cap solució}
  \task $\Delta=5$: \rest{dues solucions}
  \task $\Delta=16$: \rest{dues solucions}
\end{tasks}

\ex{2} \begin{tasks}(1)
  \task $x^{2}-x-6=0\ \Rightarrow\ \Delta=25\ \Rightarrow\ x=\tfrac{1\pm5}{2}\ \Rightarrow\ \res{x=-2,\ x=3}$
  \task $\Delta=0\ \Rightarrow\ x=\tfrac{6}{18}\ \Rightarrow\ \res{x=\tfrac{1}{3}\ (\text{doble})}$
  \task $\Delta=-11\ \Rightarrow\ \res{\text{sense solució}}$
  \task $\Delta=12\ \Rightarrow\ x=\tfrac{2\pm\sqrt{12}}{2}\ \Rightarrow\ \res{x=1\pm\sqrt{3}}$
\end{tasks}

\ex{2} \begin{tasks}(1)
  \task $x^{2}+3x-4=6\ \Rightarrow\ x^{2}+3x-10=0\ \Rightarrow\ \Delta=49\ \Rightarrow\ x=\tfrac{-3\pm7}{2}\ \Rightarrow\ \res{x=-5,\ x=2}$
  \task $x^{2}+4x+4=3x+10\ \Rightarrow\ x^{2}+x-6=0\ \Rightarrow\ \Delta=25\ \Rightarrow\ x=\tfrac{-1\pm5}{2}\ \Rightarrow\ \res{x=-3,\ x=2}$
  \task $2x^{2}-x=x^{2}+2x-3\ \Rightarrow\ x^{2}-3x+3=0\ \Rightarrow\ \Delta=-3\ \Rightarrow\ \res{\text{sense solució}}$
  \task $2x^{2}-6x+1=9\ \Rightarrow\ x^{2}-3x-4=0\ \Rightarrow\ \Delta=25\ \Rightarrow\ x=\tfrac{3\pm5}{2}\ \Rightarrow\ \res{x=-1,\ x=4}$
\end{tasks}

\ex{2} \begin{tasks}(1)
  \task $\text{m.c.m.}=6:\ \ 3x^{2}-2x=8\ \Rightarrow\ 3x^{2}-2x-8=0\ \Rightarrow\ \Delta=100\ \Rightarrow\ x=\tfrac{2\pm10}{6}\ \Rightarrow\ \res{x=-\tfrac{4}{3},\ x=2}$
  \task $\text{m.c.m.}=6:\ \ 2(x^{2}-1)+3x=12\ \Rightarrow\ 2x^{2}+3x-2=12\ \Rightarrow\ 2x^{2}+3x-14=0\ \Rightarrow\ \Delta=121\ \Rightarrow\ x=\tfrac{-3\pm11}{4}\ \Rightarrow\ \res{x=-\tfrac{7}{2},\ x=2}$
\end{tasks}

\ex{3} \begin{tasks}(1)
  \task $\text{m.c.m.}=12:\ \ 3(x-1)^{2}-2(x+2)=4(x-3)\ \Rightarrow\ 3x^{2}-8x-1=4x-12\ \Rightarrow\ 3x^{2}-12x+11=0\ \Rightarrow\ \Delta=12\ \Rightarrow\ x=\tfrac{12\pm\sqrt{12}}{6}\ \Rightarrow\ \res{x=\tfrac{6\pm\sqrt{3}}{3}}$
  \task $\text{m.c.m.}=6:\ \ 3x(x+1)-2(x-1)(x+2)=6\ \Rightarrow\ x^{2}+x+4=6\ \Rightarrow\ x^{2}+x-2=0\ \Rightarrow\ \Delta=9\ \Rightarrow\ x=\tfrac{-1\pm3}{2}\ \Rightarrow\ \res{x=-2,\ x=1}$
\end{tasks}

\ex{3} \begin{tasks}(1)
  \task $\Delta=16-4k=0\Rightarrow\res{k=4}$; \ $x^2-4x+4=0\Rightarrow\res{x=2}$ (doble)
  \task $9+3k-6=0\Rightarrow\res{k=-1}$; \ $x^2-x-6=0\Rightarrow x=3,\ \res{x=-2}$
\end{tasks}

\ex{3} \begin{tasks}(2)
  \task $(x-2)(x+5)=0\Rightarrow\res{x^2+3x-10=0}$
  \task $(2x-1)(x-3)=0\Rightarrow\res{2x^2-7x+3=0}$
\end{tasks}


% =====================================================
\seccio{6}{Equacions de tercer grau (Ruffini)}

\ex{1} \begin{tasks}(2)
  \task $\res{x=2,\ x=-3,\ x=5}$
  \task $\res{x=0,\ x=-1,\ x=\tfrac32}$
  \task $\res{x=4,\ x=3,\ x=-3}$
  \task $x^2+4=0$ no té solució: \ $\res{x=-1}$
\end{tasks}

\ex{1} \begin{tasks}(1)
  \task $x\,(x^{2}-16)=0\ \Rightarrow\ x^{2}=16\ \Rightarrow\ \res{x=\pm4,\ x=0}$
  \task $x\,(x^{2}+3x-10)=0\ \Rightarrow\ \Delta=49\ \Rightarrow\ x=\tfrac{-3\pm7}{2}\ \Rightarrow\ \res{x=-5,\ x=0,\ x=2}$
  \task $x^{3}-4x^{2}=0\ \Rightarrow\ x^{2}\,(x-4)=0\ \Rightarrow\ \res{x=0\ (\text{doble}),\ x=4}$
  \task $x\,(x^{2}+x+1)=0\ \Rightarrow\ \Delta=-3\ \Rightarrow\ \res{x=0}$
\end{tasks}

\ex{2} \begin{tasks}(1)
  \task $\begin{array}{r|rrr|r} & 1 & 0 & -7 & 6\\ 1 &  & 1 & 1 & -6\\ \hline & 1 & 1 & -6 & 0\end{array}\ \Rightarrow\ (x-1)(x^{2}+x-6)=0\ \Rightarrow\ \Delta=25\ \Rightarrow\ x=\tfrac{-1\pm5}{2}\ \Rightarrow\ \res{x=-3,\ x=1,\ x=2}$
  \task $\begin{array}{r|rrr|r} & 1 & -4 & 1 & 6\\ -1 &  & -1 & 5 & -6\\ \hline & 1 & -5 & 6 & 0\end{array}\ \Rightarrow\ (x+1)(x^{2}-5x+6)=0\ \Rightarrow\ \Delta=1\ \Rightarrow\ x=\tfrac{5\pm1}{2}\ \Rightarrow\ \res{x=-1,\ x=2,\ x=3}$
  \task $\begin{array}{r|rrr|r} & 1 & 2 & -5 & -6\\ -1 &  & -1 & -1 & 6\\ \hline & 1 & 1 & -6 & 0\end{array}\ \Rightarrow\ (x+1)(x^{2}+x-6)=0\ \Rightarrow\ \Delta=25\ \Rightarrow\ x=\tfrac{-1\pm5}{2}\ \Rightarrow\ \res{x=-3,\ x=-1,\ x=2}$
  \task $\begin{array}{r|rrr|r} & 1 & -3 & -4 & 12\\ 2 &  & 2 & -2 & -12\\ \hline & 1 & -1 & -6 & 0\end{array}\ \Rightarrow\ (x-2)(x^{2}-x-6)=0\ \Rightarrow\ \Delta=25\ \Rightarrow\ x=\tfrac{1\pm5}{2}\ \Rightarrow\ \res{x=\pm2,\ x=3}$
\end{tasks}

\ex{2} \begin{tasks}(1)
  \task $\begin{array}{r|rrr|r} & 1 & -1 & -8 & 12\\ 2 &  & 2 & 2 & -12\\ \hline & 1 & 1 & -6 & 0\end{array}\ \Rightarrow\ (x-2)(x^{2}+x-6)=0\ \Rightarrow\ \Delta=25\ \Rightarrow\ x=\tfrac{-1\pm5}{2}\ \Rightarrow\ \res{x=-3,\ x=2\ (\text{doble})}$
  \task $\begin{array}{r|rrr|r} & 1 & 3 & 3 & 1\\ -1 &  & -1 & -2 & -1\\ \hline & 1 & 2 & 1 & 0\end{array}\ \Rightarrow\ (x+1)(x^{2}+2x+1)=0\ \Rightarrow\ \Delta=0\ \Rightarrow\ x=\tfrac{-2}{2}\ \Rightarrow\ \res{x=-1\ (\text{triple})}$
  \task $\begin{array}{r|rrr|r} & 1 & -2 & 1 & -2\\ 2 &  & 2 & 0 & 2\\ \hline & 1 & 0 & 1 & 0\end{array}\ \Rightarrow\ (x-2)(x^{2}+1)=0\ \Rightarrow\ x^{2}=-1\ \Rightarrow\ \res{x=2}$
  \task $\begin{array}{r|rrr|r} & 1 & 1 & 0 & -2\\ 1 &  & 1 & 2 & 2\\ \hline & 1 & 2 & 2 & 0\end{array}\ \Rightarrow\ (x-1)(x^{2}+2x+2)=0\ \Rightarrow\ \Delta=-4\ \Rightarrow\ \res{x=1}$
\end{tasks}

\ex{2} \begin{tasks}(1)
  \task $x^{3}-4x^{2}-3x+12=0\ \Rightarrow\ \begin{array}{r|rrr|r} & 1 & -4 & -3 & 12\\ 4 &  & 4 & 0 & -12\\ \hline & 1 & 0 & -3 & 0\end{array}\ \Rightarrow\ (x-4)(x^{2}-3)=0\ \Rightarrow\ x^{2}=3\ \Rightarrow\ \res{x=4,\ x=\pm\sqrt{3}}$
  \task $x^{3}-x=24\ \Rightarrow\ x^{3}-x-24=0\ \Rightarrow\ \begin{array}{r|rrr|r} & 1 & 0 & -1 & -24\\ 3 &  & 3 & 9 & 24\\ \hline & 1 & 3 & 8 & 0\end{array}\ \Rightarrow\ (x-3)(x^{2}+3x+8)=0\ \Rightarrow\ \Delta=-23\ \Rightarrow\ \res{x=3}$
\end{tasks}

\ex{3} \begin{tasks}(1)
  \task $\begin{array}{r|rrr|r} & 2 & -5 & -4 & 3\\ -1 &  & -2 & 7 & -3\\ \hline & 2 & -7 & 3 & 0\end{array}\ \Rightarrow\ (x+1)(2x^{2}-7x+3)=0\ \Rightarrow\ \Delta=25\ \Rightarrow\ x=\tfrac{7\pm5}{4}\ \Rightarrow\ \res{x=-1,\ x=\tfrac{1}{2},\ x=3}$
  \task $\begin{array}{r|rrr|r} & 3 & 4 & -5 & -2\\ 1 &  & 3 & 7 & 2\\ \hline & 3 & 7 & 2 & 0\end{array}\ \Rightarrow\ (x-1)(3x^{2}+7x+2)=0\ \Rightarrow\ \Delta=25\ \Rightarrow\ x=\tfrac{-7\pm5}{6}\ \Rightarrow\ \res{x=-2,\ x=-\tfrac{1}{3},\ x=1}$
\end{tasks}

\ex{3} \begin{tasks}(1)
  \task $\begin{array}{r|rrrr|r} & 1 & 0 & -10 & 0 & 9\\ 1 &  & 1 & 1 & -9 & -9\\ \hline & 1 & 1 & -9 & -9 & 0\end{array}\ \Rightarrow\ (x-1)(x^{3}+x^{2}-9x-9)=0\ \Rightarrow\ \begin{array}{r|rrr|r} & 1 & 1 & -9 & -9\\ -1 &  & -1 & 0 & 9\\ \hline & 1 & 0 & -9 & 0\end{array}\ \Rightarrow\ (x+1)(x^{2}-9)=0\ \Rightarrow\ x^{2}=9\ \Rightarrow\ \res{x=\pm3,\ x=\pm1}$
  \task $\begin{array}{r|rrrr|r} & 1 & -1 & -7 & 1 & 6\\ 1 &  & 1 & 0 & -7 & -6\\ \hline & 1 & 0 & -7 & -6 & 0\end{array}\ \Rightarrow\ (x-1)(x^{3}-7x-6)=0\ \Rightarrow\ \begin{array}{r|rrr|r} & 1 & 0 & -7 & -6\\ -1 &  & -1 & 1 & 6\\ \hline & 1 & -1 & -6 & 0\end{array}\ \Rightarrow\ (x+1)(x^{2}-x-6)=0\ \Rightarrow\ \Delta=25\ \Rightarrow\ x=\tfrac{1\pm5}{2}\ \Rightarrow\ \res{x=-2,\ x=\pm1,\ x=3}$
\end{tasks}

\ex{3} $8-4k+4+4=0\Rightarrow\res{k=4}$. \ $\begin{array}{r|rrr|r} & 1 & -4 & 2 & 4\\ 2 & & 2 & -4 & -4\\ \hline & 1 & -2 & -2 & 0\end{array}$ \ $x^2-2x-2=0\Rightarrow\res{x=1\pm\sqrt3}$


% =====================================================
\seccio{7}{Problemes amb equacions}

\ex{1} $x+(x+1)=75\Rightarrow 2x=74\Rightarrow x=37$: \ \rest{37 i 38}

\ex{1} $x+3x=56\Rightarrow 4x=56\Rightarrow\res{x=14}$

\ex{1} $x-\tfrac{x}{3}-\tfrac{x}{4}=25\xrightarrow{\cdot12}12x-4x-3x=300\Rightarrow 5x=300\Rightarrow\res{60\text{ €}}$

\ex{1} $40+x=2(12+x)\Rightarrow 40+x=24+2x\Rightarrow\res{x=16\text{ anys}}$

\ex{2} $2x+2(2x+3)=54\Rightarrow 6x+6=54\Rightarrow x=8$: \ \rest{8 m i 19 m}

\ex{2} Gallines $x$: \ $2x+4(40-x)=110\Rightarrow -2x=-50\Rightarrow x=25$: \ \rest{25 gallines i 15 conills}

\ex{2} Monedes de 2 €: $x$. \ $2x+(30-x)=46\Rightarrow\res{x=16}$: \ \rest{16 de 2 € i 14 d'1 €}

\ex{2} $3x+12=2(x+12)\Rightarrow\res{x=12}$: \ fill \rest{12 anys}, mare \rest{36 anys}

\ex{2} $x(x+1)=156\Rightarrow x^2+x-156=0\Rightarrow x=\tfrac{-1\pm25}{2}$: \ $x=12$ ($-13$ no és natural): \ \rest{12 i 13}

\ex{2} $x^2-3x-40=0\Rightarrow x=\tfrac{3\pm13}{2}$: \ \rest{$x=8$ o $x=-5$}

\ex{2} $(x+3)^2-x^2=51\Rightarrow 6x+9=51\Rightarrow\res{x=7\text{ cm}}$

\ex{3} $x^2+(x+2)^2=100\Rightarrow x^2+2x-48=0\Rightarrow x=\tfrac{-2\pm14}{2}$: \ $x=6$ ($-8$ no val): \ \rest{6 cm i 8 cm}

\ex{3} Si un costat fa $x$, l'altre fa $17-x$: \ $x(17-x)=60\Rightarrow x^2-17x+60=0\Rightarrow x=\tfrac{17\pm7}{2}$: \ \rest{12 m i 5 m}

\ex{3} $12x+8(20-x)=9{,}5\cdot20\Rightarrow 4x+160=190\Rightarrow x=7{,}5$: \ \rest{7,5 kg de 12 € i 12,5 kg de 8 €}

\ex{3} $20t+25t=90\Rightarrow\res{t=2\text{ h}}$; \ a $20\cdot2=\res{40\text{ km}}$

\ex{3} $x^2(x+3)=112\Rightarrow x^3+3x^2-112=0$. \ $\begin{array}{r|rrr|r} & 1 & 3 & 0 & -112\\ 4 & & 4 & 28 & 112\\ \hline & 1 & 7 & 28 & 0\end{array}$ \ ($x^2+7x+28=0$ no té solució): \ \rest{base de 4 cm × 4 cm i altura 7 cm}

\ex{3} $x(x+1)(x+2)=8(3x+3)\Rightarrow x^3+3x^2-22x-24=0\Rightarrow (x+1)(x^2+2x-24)=0$: \ $x=-1,\ 4,\ -6$; només $x=4$ és natural: \ \rest{4, 5 i 6}


\end{document}
